%=============================================================================
% EECE 7205: Fundamentals of Computer Engineering | Fall 2026
% Dr. Fatema Nafa | Final Project Proposal | MyungJin Suh
%=============================================================================
\documentclass[10pt]{article}

\usepackage[T1]{fontenc}
\usepackage[margin=0.7in]{geometry}
\usepackage{amsmath, amssymb}
\usepackage{float}
\usepackage{xcolor}
\usepackage{listings}
\usepackage{booktabs}
\usepackage{tabularx}
\usepackage{array}
\usepackage{algorithm}
\usepackage{algpseudocode}
\usepackage[hidelinks]{hyperref}
\usepackage{fancyhdr}
\usepackage{enumitem}
\usepackage{titlesec}

\definecolor{shellbg}{HTML}{F5F5F5}
\definecolor{rulegray}{HTML}{BBBBBB}
\definecolor{headc}{HTML}{1F3864}

\lstdefinestyle{term}{
  basicstyle=\ttfamily\scriptsize,
  backgroundcolor=\color{shellbg},
  frame=single, rulecolor=\color{rulegray},
  breaklines=true, columns=fullflexible, keepspaces=true,
  showstringspaces=false, xleftmargin=4pt, xrightmargin=4pt,
  aboveskip=3pt, belowskip=3pt,
}

\setlength{\parindent}{0pt}
\setlength{\parskip}{0.28em plus 0.04em minus 0.04em}
\setlength{\textfloatsep}{6pt plus 1pt minus 1pt}
\setlength{\intextsep}{6pt plus 1pt minus 1pt}
\setlist[itemize]{leftmargin=1.15em, itemsep=0.05em, topsep=0.15em, parsep=0pt}
\setlist[enumerate]{leftmargin=1.25em, itemsep=0.04em, topsep=0.12em, parsep=0pt}
\titlespacing*{\section}{0pt}{0.65em}{0.25em}
\titlespacing*{\subsection}{0pt}{0.45em}{0.12em}

\algrenewcommand\algorithmicrequire{\textbf{Input:}}
\algrenewcommand\algorithmicensure{\textbf{Output:}}

\pagestyle{fancy}
\fancyhf{}
\fancyhead[L]{\scriptsize\color{headc} EECE 7205 \textbar\ Final Project Proposal}
\fancyhead[R]{\scriptsize\color{headc} MyungJin Suh}
\fancyfoot[C]{\scriptsize\thepage}
\renewcommand{\headrulewidth}{0.3pt}
\newcolumntype{Y}{>{\raggedright\arraybackslash}X}

\begin{document}

{\large\bfseries EECE 7205 --- Final Project Proposal}\\[0.15em]
{\bfseries Academic Advising and Course Planning Tool}\\[0.2em]
MyungJin Suh \quad\textbar\quad Dr.\ Fatema Nafa \quad\textbar\quad Fall 2026 \quad\textbar\quad October 26, 2026

\vspace{0.25em}
Individual project, one member. Topic from the Iteration~1 recommended list.
A menu-driven C++ program builds a prerequisite graph for the Northeastern
BSCmpE and writes a semester plan. It compares topological earliest-semester
packing with backtracking.

\section{Project Title}

\textbf{Academic Advising and Course Planning Tool.}
Plans the remaining required courses for one BSCmpE student under prerequisite,
term, and credit-cap constraints, and compares the two planners above.

\section{Problem Statement and Proposed Solution}

\textbf{Problem.} An advisor has to tell which requirements are done, which
courses an OR-prerequisite and a ``may be taken concurrently'' note actually
unlock, and which unlocked course should take the open seat when the cap, the
offering term, and a preferred elective cannot all win. A local choice can add
a semester.

\textbf{User.} An advisor, or a student preparing for that meeting. One degree,
one synthetic student at a time. The program does not register anyone and does
not store real names.

\textbf{Why it matters.} A plan that breaks a prerequisite cannot be followed.
A prerequisite-correct plan can still waste a fall-only course. Those are
different failures, so the project measures two algorithms.

\textbf{What the program does.} It loads the requirements, builds the clause
graph, and stops if the edges cycle. Completed courses come off the graph.
Both planners run. A separate \texttt{verifyPlan} checks strict prerequisites
in an earlier semester, concurrent ones in the same semester or earlier, the
offering term, and the credit cap. The screen shows each plan, its span (index
of the last semester used, holes included), a preference penalty, and the time.
Scope is the BSCmpE required list plus prerequisites those entries name, not
the whole university, and not seat counts.

\textbf{Features.}
\begin{enumerate}
  \item Load requirements and clauses transcribed from the catalog.
  \item Load one synthetic student: completed courses, credit cap, preferred electives.
  \item Build the graph and refuse a cycle.
  \item Look up a course by code.
  \item Plan by clause-count topological order and earliest-semester packing.
  \item Plan by backtracking, pruning on the best span.
  \item Verify both plans and print span, preference, and time.
  \item Save the chosen plan.
  \item Time the worst-case family on synthetic graphs.
\end{enumerate}

\textbf{Inputs, outputs, and sources.}
Requirements and prerequisite text, from the Northeastern undergraduate catalog:

{\scriptsize
\href{https://catalog.northeastern.edu/undergraduate/engineering/electrical-computer/computer-engineering-bscompe/}{https://catalog.northeastern.edu/undergraduate/engineering/}\\
\href{https://catalog.northeastern.edu/undergraduate/engineering/electrical-computer/computer-engineering-bscompe/}{electrical-computer/computer-engineering-bscompe/}\\
\url{https://catalog.northeastern.edu/course-descriptions/eece/}}
Three clauses are copied exactly because the tests use them.
EECE 2160 (4 hours): EECE 2140 or CS 2510, either may be concurrent.
EECE 2322 (4 hours): EECE 2160 or EECE 2310.
EECE 2560 (4 hours): EECE 2160 or CS 1500. The file will hold every
required course on the BSCmpE page and the prerequisites those rows name.
Math, physics, and CS appear only when an EECE row lists them. Rows are
transcribed, not scraped. If a catalog page does not state a term, the course
is marked \texttt{FS} and the README says so.

Clauses are separated by \texttt{;}, AND-members by \texttt{+}, and a trailing
\texttt{@} means concurrent. \texttt{offered} is \texttt{F}, \texttt{S}, or
\texttt{FS}. A course unlocks when one clause is satisfied. A strict member
must be earlier. A concurrent member may be earlier or in the same semester.
Storing every alternative as a required edge would turn the OR into an AND.

\begin{lstlisting}[style=term]
code,title,credits,offered,clauses
EECE2560,Fundamentals of Engineering Algorithms,4,FS,EECE2160;CS1500

id,completed,maxCredits,prefer
S01,EECE2140,8,EECE2560
\end{lstlisting}

Student ids are \texttt{S01}, \texttt{S02}, \ldots, synthetic, not people.
Saving writes \texttt{semester,code,title,credits}. Timing appends algorithm,
\(V\), median time, and whether the trial hit the cap.

\textbf{Why this problem.} It is the recommended advising topic, limited to one
degree so one person can transcribe the data and still measure both planners.
The work is a graph, a topological order, a backtracking search, and a
best-case and worst-case bound for each.

\section{Team Structure, Contributions, and Objectives}

\textbf{Team.} One member. The Iteration~1 sheet should use this same topic
and description: MyungJin Suh; Academic Advising and Course Planning Tool; a
C++ planner for the BSCmpE prerequisite graph, catalog rows transcribed from
the undergraduate catalog, synthetic students, topological packing compared
with backtracking, best- and worst-case time and space.

\textbf{Objectives.} A working planner under prerequisite, term, and cap
constraints. The graph, both planners, the queue, and insertion sort are
hand-written; a \texttt{vector} is only a resizable array while reading.
Compare span, preference, and time. State best and worst time, and space, as
\(\Theta\) of the input that attains the bound. Test normal, boundary, and
invalid cases, including a planted cycle. Use feature branches on
\texttt{EECE7205\_HW}.

\textbf{Contributions --- this iteration.} Done: chose the topic and the
BSCmpE bound; cited the catalog and checked the three clauses; defined the CSV,
the OR rule, and the concurrent rule; specified both planners, the verifier,
and the bounds; wrote the pseudocode, tests T1--T10, and the timing plan,
including one input where the planners must disagree. Next, still mine:
transcribe \texttt{data/catalog.csv}; implement load, the BST index, and cycle
detection; implement both planners and the verifier and confirm T4; add the
timing script. This iteration is the design. Programming and testing are the
following iterations.

\textbf{Learning objectives.} Pass records by reference and const reference.
Index codes with a BST and state the chain worst case. Implement an OR-clause
topological order and a pruned backtracking search. Read and write CSV and
reject a bad row. Derive the \(\Theta\) bounds, then time the worst-case
family. Use branches, small commits, and pull requests.

\section{Program Design and GitHub Collaboration --- This Iteration}

\subsection{Initial program design}

The nine features above are the menu, choices 1--9.

\textbf{Data structures.} \texttt{Course} holds code, title, credits, fall,
spring, and a required flag. \texttt{Clause} holds the owner index, up to four
member indexes, a concurrent flag per member, and \texttt{left}, the members
not yet completed. \texttt{Plan} holds \texttt{termOf[MAXC]} (\(-1\) if
unplaced), span, and preference (sum of term indexes of preferred electives).
Catalog and clauses live in fixed arrays; overflow is an error. The catalog is
a few dozen required courses.

A BST keyed by code maps to a catalog index. The node is the in-class shape:
key, index, left, right. Lookup is \(\Theta(1)\) at the root and \(\Theta(V)\)
on a chain. Sorted insertion builds that chain. The tree is not balanced.
Space \(\Theta(V)\). The graph is the clause array plus
\texttt{int adj[MAXC][MAXD]}, the clauses that mention each course. Clause
counts, not one indegree, keep an OR from becoming an AND. Space
\(\Theta(V+P)\), \(P\) the number of clause memberships. The Kahn queue is a
ring buffer of indexes, not \texttt{std::queue}. A hand-written insertion sort
orders a ready list by code: \(\Theta(r)\) best, \(\Theta(r^{2})\) worst,
\(r \le V\).

\textbf{Planners.} Both place only unlocked courses and both refuse a cycle.
\texttt{verifyPlan} does not share code with either search. Planner~A builds
the clause-count order (a course enters when some clause has \texttt{left = 0},
ties break toward the smaller code), then places each course in the earliest
legal semester. That earliest fit is the greedy step, and T4 is where it picks
the wrong fall-only course. Planner~B tries semesters in that same order,
earliest first, keeps the smallest span, and cuts a branch that cannot finish
sooner. Preference breaks equal spans.

\textbf{Bounds.}
\label{sec:bounds}
\(V\) courses left to place, \(P\) clause memberships, \(S\) semesters
(\(S = 16\) in the program, eight years of fall and spring), \(r\) courses
ready in one semester.

{\footnotesize
\begin{tabularx}{\textwidth}{@{}>{\raggedright\arraybackslash}p{1.35in} Y Y@{}}
\toprule
\textbf{Operation} & \textbf{Time} & \textbf{Space} \\
\midrule
Load and build &
\(\Theta(V+P)\) best and worst &
\(\Theta(V+P)\) \\
BST lookup &
Best \(\Theta(1)\), the root. Worst \(\Theta(V)\), a chain &
\(\Theta(V)\) \\
Topological order &
\(\Theta(V+P)\) best and worst. Each course is enqueued once; each membership is decremented once &
\(\Theta(V)\) \\
Earliest-semester packing &
\(\Theta(VS+P)\) best and worst &
\(\Theta(V)\) \\
Backtracking &
Best \(\Theta(V+P)\): a chain with one legal semester per course, and the span bound cuts the rest.
Worst \(\Theta(S^{V})\): \(P = 0\), every course offered every term, cap never binds &
\(\Theta(V)\) \\
\bottomrule
\end{tabularx}}

Packing is polynomial. Backtracking is exponential on that family, which is why
the experiment has a time cap. Extra space beyond the graph is linear in \(V\).

\textbf{Functions.}
\texttt{loadCatalog} reads rows and rejects a bad credit, term, or unknown
clause member.
\texttt{loadStudent} reads one synthetic student.
\texttt{buildIndex} / \texttt{searchNode} maintain the BST; a duplicate code is
an error.
\texttt{buildGraph} fills counts and adjacency lists.
\texttt{topologicalOrder} returns false if a required course never becomes ready.
\texttt{greedyPlan} packs along that order.
\texttt{backtrackPlan} is Algorithm~\ref{alg:bt}.
\texttt{unlocked} applies the concurrent rule.
\texttt{verifyPlan} rechecks prerequisites, offerings, and the cap.
\texttt{scorePlan} returns span, then preference.
\texttt{insertionSort} orders a ready list.
\texttt{runTimingSeries} takes five trials and the median, and stops at a time cap.
\texttt{showMenu}, \texttt{savePlan}, and \texttt{freeTree} do what the names say.
Planners do not read \texttt{cin}. Tests call them directly.

\textbf{Pseudocode, written before the C++.}

\begin{algorithm}[H]
\caption{Main program}
\footnotesize
\begin{algorithmic}[1]
\Require catalog path, and a student path when the user loads one
\Ensure both plans printed, or the bad row named
\State $courses, clauses \gets \Call{loadCatalog}{path}$
\If{the load failed, or a code is duplicated}
    \State print the row and stop
\EndIf
\State \Call{buildGraph}{courses, clauses}
\Repeat
    \State $choice \gets \Call{showMenu}{}$
    \If{choice is outside 1--9} \State print ``Enter a number from 1 to 9''
    \ElsIf{choice is 4, 5, or 6}
        \If{no student is loaded} \State print ``Load a student first''
        \Else \State run the requested planners, \Call{verifyPlan}{result}, print span, preference, and time
        \EndIf
    \ElsIf{choice is 9} \State \Call{freeTree}{root} and stop
    \EndIf
\Until{stopped}
\end{algorithmic}
\end{algorithm}

\begin{algorithm}[H]
\caption{Earliest-semester packing}
\footnotesize
\begin{algorithmic}[1]
\Require courses, clauses, completed set, cap, and $S$ semesters
\Ensure a plan, or failure if the order cannot be built
\State $order \gets \Call{topologicalOrder}{courses, completed}$
\If{order is incomplete} \State \Return failure \EndIf
\State $load[1 \ldots S] \gets 0$
\For{each not-yet-completed course $c$ along order}
    \State $placed \gets$ false
    \For{$s \gets 1$ to $S$}
        \If{$c$ is offered in $s$, the cap has room, and some clause is satisfied
            (strict members earlier; concurrent members in $s$ or earlier)}
            \State put $c$ in $s$, add its credits to $load[s]$, set $placed$, and \textbf{break}
        \EndIf
    \EndFor
    \If{placed is false} \State \Return failure \EndIf
\EndFor
\State \Return the plan
\end{algorithmic}
\end{algorithm}

\begin{algorithm}[H]
\caption{Backtracking plan}
\label{alg:bt}
\footnotesize
\begin{algorithmic}[1]
\Require the topological order, the cap, and $S$
\Ensure the stored plan with the smallest span
\State $best.span \gets S + 1$; \quad \Call{Search}{1}; \quad \Return $best$
\Statex
\Function{Search}{$i$}
    \If{$i$ is past the end of order}
        \If{this span is smaller than $best.span$} \State $best \gets$ this plan \EndIf
        \State \Return
    \EndIf
    \State $c \gets order[i]$
    \If{$c$ is already completed} \State \Call{Search}{$i + 1$}; \quad \Return \EndIf
    \For{$s \gets 1$ to $S$}
        \If{$s$ cannot beat $best.span$} \State \Return \EndIf
        \If{$c$ is offered in $s$, the cap holds, and some clause is satisfied}
            \State place $c$ in $s$; \quad \Call{Search}{$i + 1$}; \quad undo
        \EndIf
    \EndFor
\EndFunction
\end{algorithmic}
\end{algorithm}

Semesters are tried earliest first, so a later \(s\) cannot beat a span already
found and the inner return is the prune. On a chain with one legal semester
per course, the first plan is the only plan. That is the \(\Theta(V+P)\) case.

\textbf{Tests.}
\label{sec:tests}
T1--T2 use the published clauses. T3--T6 and T9 use a constructed file, not the
catalog: \(P\) and \(Q\) are fall-only, \(R\) is spring-only, \(Q\) is a strict
prerequisite of \(R\), the cap is one course, and the calendar starts in a fall.

{\scriptsize
\begin{tabularx}{\linewidth}{@{}l l Y Y@{}}
\toprule
\textbf{ID} & \textbf{Kind} & \textbf{Input} & \textbf{Expected} \\
\midrule
T1 & Normal & Look up EECE 2560 &
Clause EECE 2160 or CS 1500. Credits 4. \\
T2 & Normal & Completed set is EECE 2160 &
EECE 2560 and EECE 2322 unlock. EECE 2160 is not planned again. \\
T3 & Normal & \(P, Q, R\), cap 1, both planners &
Both pass \texttt{verifyPlan}. \(R\) is never beside or before \(Q\). \\
T4 & Normal & Same file, compare spans &
Greedy span 4: \(P\), empty spring, \(Q\), \(R\).
Backtracking span 3: \(Q\), \(R\), \(P\). \\
T5 & Boundary & Every required course already done &
Both plans empty, span 0. \\
T6 & Boundary & Cap 0, one course left &
Both fail. No plan file is written. \\
T7 & Invalid & Duplicate code EECE 2560 &
Load stops on that row. \\
T8 & Invalid & Clause names EECE 9999 &
Load stops and names EECE 9999. \\
T9 & Invalid & Add edge \(R \to Q\) &
Order is incomplete. Neither planner prints a semester list. \\
T10 & Invalid & Menu 0, 10, and a blank line &
``Enter a number from 1 to 9.'' \\
\bottomrule
\end{tabularx}}

\textbf{Experiment.} Worst-case family from the bound: \(P = 0\), every course
offered every term, cap loose, \(S = 4\), \(V \in \{8, 10, 12\}\), five trials,
median time. A trial past the cap is recorded as capped. Packing should stay
fast; backtracking's median should jump. No millisecond guess before measuring.
A second run uses the transcribed BSCmpE file and three synthetic students
(nothing completed, the first programming course completed, one preferred
elective). Both plans must pass \texttt{verifyPlan}. Which span is smaller
waits on that file.

\subsection{GitHub setup}

Repository \texttt{EECE7205\_HW}, private, only collaborator, default branch
\texttt{main}: \url{https://github.com/mjsuhhh/EECE7205_HW}.
Homework stays where it is. Project files go under \texttt{FinalProject/}.
Instructor read access if the course asks.

\texttt{main} stays compilable. Feature branches:
\texttt{feature/load-catalog}, \texttt{feature/prereq-graph},
\texttt{feature/greedy-plan}, \texttt{feature/backtracking}.
Each opens a pull request. I review my own diff before merging: builds with
\texttt{g++ -std=c++17 -Wall -Wextra}, new tests pass, the function matches the
pseudocode, and \texttt{verifyPlan} is not the planner it checks.
Commits stay small (\texttt{Add the BSCmpE catalog CSV},
\texttt{Add clause-count topological order}).
Order of work: catalog and README, then the graph, packing, backtracking, then
the timing script.

\texttt{FinalProject/README.md} comes with the first code: description, both
catalog URLs, CSV columns, OR and concurrent rules, build and run, and progress.
Build and run: \texttt{g++ -std=c++17 -Wall -Wextra -o planner src/main.cpp}
and \texttt{./planner data/catalog.csv}.
Progress now: topic, sources, pseudocode, bounds, and tests. No \texttt{.cpp}
yet.

\section{Tools and Programming Language}

No external libraries. The graph, both planners, the queue, the BST, and
insertion sort are mine. Files use \texttt{<fstream>} and \texttt{<sstream>}.
Timing uses \texttt{<chrono>}. A CSV or graph library would hide the clauses
and the search.

{\footnotesize
\begin{tabularx}{\textwidth}{@{}>{\raggedright\arraybackslash}p{1.15in} Y >{\raggedright\arraybackslash}p{1.45in}@{}}
\toprule
\textbf{Tool} & \textbf{Use} & \textbf{Skill} \\
\midrule
C++17 & Structs, functions, references, arrays, pointers. No class hierarchy. &
Existing, on a larger program \\
g++ and clang & \texttt{-std=c++17 -Wall -Wextra}, as in the homework. & Existing \\
Visual Studio Code & Edit and read compiler diagnostics. & Existing \\
Git and GitHub & Branches, commits, and pull requests on \texttt{EECE7205\_HW}. &
Commits existing; pull requests new \\
GDB or the VS Code debugger & Stop in \texttt{unlocked} and \texttt{Search}. & To develop \\
Catalog and student CSV & Real clauses transcribed. Synthetic student rows. & To develop \\
T1--T10 and timing & Correctness, then \(P = 0\) at \(V = 8, 10, 12\). &
Bounds existing; timing a search is new \\
\texttt{run\_tests.sh} & Compile and run the non-menu checks, as in Assignment~01. & Existing pattern \\
GitHub Actions & Not used. The shell script is the check. & Not used \\
\bottomrule
\end{tabularx}}

\end{document}
